\documentclass[12pt,a4paper]{article}

\usepackage{graphicx}
\usepackage{grffile}
\usepackage{geometry}
\geometry{margin=2.5cm}
\usepackage{amsmath}
\usepackage{booktabs}
\usepackage{longtable}
\usepackage{float}
\usepackage{caption}
\usepackage{enumitem}
\usepackage{xcolor}
\usepackage{fancyhdr}
\usepackage{titlesec}
\usepackage[hidelinks]{hyperref}


\usepackage{xepersian}
\settextfont{XB Niloofar}


\newcommand{\en}[1]{\lr{#1}}

\graphicspath{{pics/}}

\newcommand{\pulsepic}[2]{
  \begin{minipage}[t]{0.42\linewidth}
    \centering
    \includegraphics[width=\linewidth]{#1}\\[4pt]
    {\small پالس~#2}
  \end{minipage}
}

\title{\textbf{گزارش آزمایش چهارم آزمایشگاه مدار منطقی\\ مدار کنترل‌کننده (طراحی به روش \en{ASM Chart})}}
\author{}
\date{}

\setlength{\headheight}{18pt}
\pagestyle{fancy}
\fancyhf{}
\fancyhead[C]{گزارش آزمایش چهارم --- مدار کنترل‌کننده}
\fancyfoot[C]{\thepage}

\begin{document}

\begin{titlepage}
\centering
\vspace*{2cm}
{\Huge \textbf{گزارش کار آزمایشگاه}}\\[0.5cm]
{\Large مدار منطقی}\\[2cm]
{\LARGE \textbf{آزمایش چهارم: مدار کنترل‌کننده}}\\[0.3cm]
{\large طراحی و پیاده‌سازی با کمک نمودار حالت \en{(ASM Chart)} در \en{Proteus}}\\[2cm]

{\large علی‌اصغر طایری - ۴۰۴۱۷۱۱۳۳}\\[0.3cm]
{\large رضا حضرتی - ۴۰۴۱۰۵۷۶۷}\\[2cm]

\vfill

\end{titlepage}

\tableofcontents
\newpage

\section{مقدمه و تعریف مسئله کلی آزمایش}

هدف از آزمایش چهارم، آشنایی عملی با روند طراحی مدارهای کنترل‌کننده (ماشین‌های حالت متناهی) با استفاده از نمودار \en{ASM Chart} و پیاده‌سازی آن با گیت‌های منطقی، فلیپ‌فلاپ و مدارهای \en{MSI} در محیط شبیه‌سازی \en{Proteus} است. در این روش، ابتدا رفتار سیستم به‌صورت یک نمودار حالت (شامل حالت‌ها، شرایط انتقال بین حالت‌ها و خروجی هر حالت) مدل‌سازی می‌شود و سپس این نمودار به یک مدار ترتیبی همزمان (متشکل از رجیستر حالت و منطق ترکیبی تعیین حالت بعد و خروجی) تبدیل می‌گردد.

در این آزمایش دو مسئله‌ی مستقل با همین روش پیاده‌سازی شده‌اند:

\begin{itemize}[rightmargin=1cm]
\item \textbf{بخش اول (\en{5-1})}: طراحی تایمر یک ماشین لباس‌شویی که با گرفتن ورودی‌های \en{Start}، \en{Valve}، \en{Door} و \en{Function} به‌ترتیب عملیات آب‌گیری، (گرم‌کردن آب در صورت انتخاب برنامه‌ی گرم)، شست‌وشو، تخلیه و خشک‌کردن را در بازه‌های زمانی مشخص انجام می‌دهد و در پایان با فعال‌کردن خروجی \en{Finish} منتظر فشرده‌شدن کلید \en{Reset} می‌ماند.
\item \textbf{بخش دوم (\en{5-2})}: طراحی مدار کنترل‌کننده‌ی یک تلفن سکه‌ای (راه‌دور) که با شمارش سکه‌های ورودی، برقراری و قطع خودکار تماس بر اساس موجودی سکه‌ها و کسر دوره‌ای هزینه‌ی مکالمه را مدیریت می‌کند.
\end{itemize}

هر دو مدار طبق دستور آزمایش، ابتدا با ترسیم نمودار \en{ASM} طراحی و سپس در نرم‌افزار \en{Proteus} پیاده‌سازی و با چند سناریوی آزمایشی صحت‌سنجی شده‌اند. در ادامه‌ی گزارش، هر بخش به‌طور مجزا و کامل (از تعریف مسئله تا نتیجه‌گیری) شرح داده شده است.


\section{بخش اول: تایمر ماشین لباس‌شویی (\lr{5-1})}


\subsection{تعریف مسئله}

هدف، طراحی یک تایمر برای ماشین لباس‌شویی است. ورودی‌های مدار عبارت‌اند از دو کلید دوحالته \en{Start} (شروع)، \en{Valve} (وضعیت باز/بسته بودن شیر آب) و \en{Door} (وضعیت باز/بسته بودن در دستگاه)، یک کلید دوحالته \en{Function} برای انتخاب برنامه‌ی شست‌وشوی گرم یا سرد، یک کلید فشاری \en{Reset} برای بازگشت به حالت اولیه و یک مولد پالس برای ورودی ساعت. خروجی‌های مدار نیز شامل \en{Fill} (آب‌گیری)، \en{Heat} (گرم‌کردن آب)، \en{Wash} (شست‌وشو)، \en{Drain} (تخلیه)، \en{Dry} (خشک‌کردن) و \en{Finish} (خاتمه) هستند.

طبق صورت مسئله، با فشردن کلید \en{Start} --- مشروط بر باز بودن شیر آب، بسته بودن در دستگاه و مشخص بودن برنامه‌ی شست‌وشو --- کار ماشین آغاز می‌شود. در برنامه‌ی سرد، عملیات آب‌گیری، شست‌وشو، تخلیه و خشک‌کردن به‌ترتیب در زمان‌های \lr{$T_1,T_3,T_4,T_5$} انجام می‌شود و در برنامه‌ی گرم، پیش از شست‌وشو یک مرحله‌ی گرم‌کردن آب به مدت \lr{$T_2$} نیز اضافه می‌شود. در پایان هر دو برنامه، خروجی \en{Finish} فعال شده و مدار تا فشرده‌شدن کلید \en{Reset} در همان وضعیت باقی می‌ماند. طبق صورت مسئله \lr{$T_1=T_4=T_5=2$} پالس ساعت و \lr{$T_2=T_3=3$} پالس ساعت در نظر گرفته شده است.

\subsection{بررسی و انتخاب روش طراحی}

با توجه به توالی نسبتاً خطی و ساده‌ی عملیات (آب‌گیری، گرم‌کردن اختیاری، شست‌وشو، تخلیه، خشک‌کردن، خاتمه)، تعداد حالت‌های لازم برای این ماشین محدود و برابر با هفت حالت است. به همین دلیل، برخلاف بخش دوم (که به علت پیچیدگی معادلات از روش \en{one-hot} استفاده شد)، در این بخش از \textbf{کدگذاری حالت به روش باینری} استفاده شد: یک رجیستر حالت ۴ بیتی (چهار فلیپ‌فلاپ نوع \en{D}) که ظرفیت نمایش تا ۱۶ حالت را دارد، برای نگه‌داری شماره‌ی حالت جاری به کار رفت و خروجی این رجیستر توسط یک \textbf{دیکودر ۴ به ۱۶} (\en{Decoder 4 to 16}) به ۱۶ خط جداگانه (یک خط فعال به ازای هر حالت) تبدیل شد.

این روش مزیت هر دو رویکرد را ترکیب می‌کند: از یک‌سو تعداد فلیپ‌فلاپ‌های لازم بسیار کمتر از روش \en{one-hot} است (۴ فلیپ‌فلاپ به‌جای ۷ فلیپ‌فلاپ)، و از سوی دیگر، چون خروجی دیکودر در عمل مانند یک بردار \en{one-hot} رفتار می‌کند، معادلات خروجی و معادلات حالت بعد به‌سادگی به‌صورت \en{OR} چند خط دیکد‌شده نوشته می‌شوند و پیچیدگی معمول کدگذاری باینری خالص (که در آن هر خروجی تابعی از ترکیب چند بیت حالت است) وجود ندارد.

\subsection{بلوک‌دیاگرام طرح پیشنهادی اولیه}

بلوک‌دیاگرام کلی مدار از سه بخش اصلی تشکیل می‌شود: (۱) منطق تعیین حالت بعد که ورودی‌های \en{Start}، \en{Valve}، \en{Door}، \en{Function} و \en{Reset} را همراه با حالت جاری دریافت کرده و مقدار بعدیِ رجیستر حالت را می‌سازد، (۲) رجیستر حالت (۴ فلیپ‌فلاپ \en{D}) به‌همراه دیکودر ۴ به ۱۶، و (۳) منطق خروجی که با \en{OR} کردن خطوط دیکد‌شده‌ی مربوط به هر حالت، خروجی‌های \en{Fill}، \en{Heat}، \en{Wash}، \en{Drain} و \en{Dry} را تولید می‌کند.

\begin{figure}[H]
\centering
\includegraphics[width=0.85\linewidth]{4-1-ASMChart.PNG}
\caption{نمودار  \lr{ASM} طراحی‌شده برای تایمر ماشین لباس‌شویی --- مرجع طراحی حالت بعد و خروجی هر حالت (برای کشیدن شکل داده های چارت به هوش مصنوعی داده شده تا شکل غیر دست نویس و با خوانایی بیشتر باشه)}
\end{figure}

\subsection{روند پیاده‌سازی}

ابتدا برای هر یک از هفت حالت مدار --- $S_0$: بی‌کاری (\en{Idle})، $S_1$: آب‌گیری (\en{Fill})، $S_2$: گرم‌کردن آب (\en{Heat}، فقط در برنامه‌ی گرم)، $S_3$: شست‌وشو (\en{Wash})، $S_4$: تخلیه (\en{Drain})، $S_5$: خشک‌کردن (\en{Dry}) و $S_6$: خاتمه (\en{Finish}) --- نمودار \en{ASM} رسم و شرط خروج از هر حالت مشخص شد. سپس مراحل زیر طی شد:

\begin{enumerate}[rightmargin=1cm]
\item چهار کلید دوحالته‌ی \en{Start}، \en{Valve}، \en{Door} و \en{Function} و یک کلید فشاری \en{Reset} روی صفحه‌ی شماتیک قرار گرفت و برای اطمینان از سطح منطقی مشخص در هر وضعیت، هر کلید با یک مقاومت \lr{$10\,\text{k}\Omega$} به تغذیه/زمین متصل شد (شبیه‌سازی ورودی‌های \en{pull-up}).
\item شمارش زمان هر حالت (رسیدن به \lr{$T_1$} تا \lr{$T_5$} پالس ساعت) با استفاده از همان رجیستر حالت انجام شد: برای حالت‌هایی که بیش از یک پالس ساعت طول می‌کشند، مسیر بازگشتی «ماندن در همان حالت» تا رسیدن شمارش داخلی به مقدار موردنظر فعال نگه داشته می‌شود و در پالس بعدی، شرط خروج از حالت برقرار شده و رجیستر به حالت بعدی منتقل می‌شود.
\item منطق حالت بعد با گیت‌های \en{AND} و \en{OR} بر اساس معادلات به‌دست‌آمده از نمودار \en{ASM} پیاده‌سازی و به ورودی‌های \en{D} چهار فلیپ‌فلاپ متصل شد.
\item خروجی این چهار فلیپ‌فلاپ به ورودی دیکودر ۴ به ۱۶ داده شد تا شانزده خط $Q_0$ تا $Q_{15}$ (هر خط معادل یک حالت) در دسترس قرار گیرد.
\item در نهایت، پنج گیت \en{OR} با گرفتن خطوط دیکد‌شده‌ی مربوط به حالت‌های فعال‌کننده‌ی هر عملکرد، خروجی‌های \en{Fill}، \en{Heat}، \en{Wash}، \en{Drain} و \en{Dry} را می‌سازند (برای نمونه، \en{Fill} از \en{OR} خط حالت \lr{$S_1$} تشکیل شده و \en{Wash} از \en{OR} خطوط مربوط به حالت شست‌وشو در هر دو برنامه به‌دست می‌آید). خروجی \en{Finish} نیز مستقیماً از خط دیکد‌شده‌ی حالت \lr{$S_6$} گرفته شده است.
\end{enumerate}

\subsection{دیاگرام/شماتیک طرح اولیه}

در طرح اولیه، همان معماری رجیستر حالت + دیکودر در نظر گرفته شد؛ تفاوت عمده با طرح نهایی صرفاً در شکل نهایی معادلات منطق حالت بعد بود که پس از اولین دور شبیه‌سازی و مشاهده‌ی رفتار نادرست در برخی گذرها (به‌ویژه گذر بین دو برنامه‌ی گرم و سرد)، اصلاح شدند. از این رو شماتیک طرح اولیه و طرح نهایی از نظر ساختار بلوکی یکسان است و تغییرات صرفاً در سطح سیم‌کشی گیت‌های منطق حالت بعد اعمال شده است.

\subsection{دیاگرام طرح نهایی}

شماتیک نهایی مدار (خروجی مستقیم از \en{Proteus}) در شکل زیر آمده است. در این شماتیک، بلوک سمت چپ ورودی‌ها و معکوس‌کننده‌ها، بلوک وسط منطق حالت بعد (گیت‌های \en{AND}/\en{OR})، فلیپ‌فلاپ‌های رجیستر حالت و دیکودر ۴ به ۱۶، و بلوک سمت راست گیت‌های \en{OR} تولیدکننده‌ی خروجی‌های نهایی را نشان می‌دهد.





\begin{figure}[H]
\centering
\includegraphics[width=\linewidth]{4-1-main.png}
\caption{شماتیک نهایی مدار تایمر ماشین لباس‌شویی در \lr{Proteus}}
\end{figure}

\subsection{بررسی \lr{Test Case} ها}

برای صحت‌سنجی مدار، سه سناریوی زیر در \en{Proteus} اجرا و رفتار مدار پله‌به‌پله بررسی شد: اجرای کامل برنامه‌ی شست‌وشوی سرد، بررسی عملکرد کلید \en{Reset}، و اجرای کامل برنامه‌ی شست‌وشوی گرم.

\subsubsection{الف) اجرای کامل برنامه‌ی شست‌وشوی سرد}

با بستن در، باز بودن شیر آب، انتخاب برنامه‌ی سرد و فشردن \en{Start}، مدار به‌ترتیب وارد حالت‌های آب‌گیری (۲ پالس)، شست‌وشو (۳ پالس)، تخلیه (۲ پالس) و خشک‌کردن (۲ پالس) شده و در پایان خروجی \en{Finish} فعال و ثابت می‌ماند. تصاویر زیر تمام ۱۱ پالس ساعت ثبت‌شده از این اجرا را به‌ترتیب، از پالس ۰ تا پالس ۱۰، نشان می‌دهند.

\begin{figure}[H]
\centering
\setlength{\tabcolsep}{12pt}
\begin{tabular}{cc}
\pulsepic{4-1-ColdTest-1.PNG}{۰} & \pulsepic{4-1-ColdTest-2.PNG}{۱} \\[10pt]
\pulsepic{4-1-ColdTest-3.PNG}{۲} & \pulsepic{4-1-ColdTest-4.PNG}{۳} \\[10pt]
\pulsepic{4-1-ColdTest-5.PNG}{۴} & \pulsepic{4-1-ColdTest-6.PNG}{۵} \\[10pt]
\multicolumn{2}{c}{\pulsepic{4-1-ColdTest-7.PNG}{۶}}
\end{tabular}
\caption{تصاویر گام‌به‌گام اجرای برنامه‌ی شست‌وشوی سرد --- پالس‌های ۰ تا ۶}
\end{figure}

\begin{figure}[H]
\centering
\setlength{\tabcolsep}{12pt}
\begin{tabular}{cc}
\pulsepic{4-1-ColdTest-8.PNG}{۷} & \pulsepic{4-1-ColdTest-9.PNG}{۸} \\[10pt]
\pulsepic{4-1-ColdTest-10.PNG}{۹} & \pulsepic{4-1-ColdTest-11.PNG}{۱۰}
\end{tabular}
\caption{تصاویر گام‌به‌گام اجرای برنامه‌ی شست‌وشوی سرد --- پالس‌های ۷ تا ۱۰ (ادامه)}
\end{figure}

\subsubsection{ب) بررسی عملکرد کلید \lr{Reset}}

\begin{figure}[H]
\centering
\includegraphics[width=0.8\linewidth]{4-1-ResetTest.PNG}
\caption{بازگشت مدار به حالت اولیه پس از فشردن کلید \lr{Reset}}
\end{figure}

\subsubsection{ج) اجرای کامل برنامه‌ی شست‌وشوی گرم}

با همان شرایط شروع اما با انتخاب برنامه‌ی گرم، مدار علاوه‌بر مراحل قبلی، پیش از شست‌وشو نیز به مدت ۳ پالس ساعت وارد حالت \en{Heat} می‌شود. تصاویر زیر تمام ۱۴ پالس ساعت ثبت‌شده از این اجرا را به‌ترتیب، از پالس ۰ تا پالس ۱۳، نشان می‌دهند.

\begin{figure}[H]
\centering
\setlength{\tabcolsep}{12pt}
\begin{tabular}{cc}
\pulsepic{4-1-HotTest-1.PNG}{۰} & \pulsepic{4-1-HotTest-2.PNG}{۱} \\[10pt]
\pulsepic{4-1-HotTest-3.PNG}{۲} & \pulsepic{4-1-HotTest-4.PNG}{۳} \\[10pt]
\pulsepic{4-1-HotTest-5.PNG}{۴} & \pulsepic{4-1-HotTest-6.PNG}{۵} \\[10pt]
\multicolumn{2}{c}{\pulsepic{4-1-HotTest-7.PNG}{۶}}
\end{tabular}
\caption{تصاویر گام‌به‌گام اجرای برنامه‌ی شست‌وشوی گرم --- پالس‌های ۰ تا ۶}
\end{figure}

\begin{figure}[H]
\centering
\setlength{\tabcolsep}{12pt}
\begin{tabular}{cc}
\pulsepic{4-1-HotTest-8.PNG}{۷} & \pulsepic{4-1-HotTest-9.PNG}{۸} \\[10pt]
\pulsepic{4-1-HotTest-10.PNG}{۹} & \pulsepic{4-1-HotTest-11.PNG}{۱۰} \\[10pt]
\pulsepic{4-1-HotTest-12.PNG}{۱۱} & \pulsepic{4-1-HotTest-13.PNG}{۱۲} \\[10pt]
\multicolumn{2}{c}{\pulsepic{4-1-HotTest-14.PNG}{۱۳}}
\end{tabular}
\caption{تصاویر گام‌به‌گام اجرای برنامه‌ی شست‌وشوی گرم --- پالس‌های ۷ تا ۱۳ (ادامه)}
\end{figure}

\subsection{نتیجه‌گیری بخش اول}

مدار تایمر ماشین لباس‌شویی با استفاده از رجیستر حالت باینری و دیکودر ۴ به ۱۶ به‌درستی طراحی و پیاده‌سازی شد و در هر سه سناریوی آزمایشی (برنامه‌ی سرد، برنامه‌ی گرم و بازگشت با \en{Reset}) رفتار مطابق با صورت مسئله از خود نشان داد.

\newpage

\section{بخش دوم: تلفن راه‌دور (\lr{5-2})}


\subsection{تعریف مسئله}

در این بخش باید مدار کنترل‌کننده‌ی یک تلفن سکه‌ای طراحی شود. ورودی‌های مدار عبارت‌اند از یک کلید دوحالته \en{Start}، یک کلید فشاری \en{Coin} (هر بار فشردن آن معادل انداختن یک سکه‌ی ده‌ریالی است)، یک کلید فشاری \en{Reset} و یک مولد پالس ساعت. خروجی‌های مدار شامل \en{Call} (نشانگر برقراری تماس)، \en{Alarm} (هشدار اتمام موجودی) و دو نمایشگر هفت‌قطعه برای نمایش تعداد سکه‌های موجود (حداکثر ۹۹ سکه) هستند.

طبق صورت مسئله، برای برقراری تماس باید حداقل یک سکه انداخته و سپس کلید \en{Start} فشرده شود؛ با برقراری تماس بلافاصله ده ریال از موجودی کسر می‌شود و از آن پس به ازای هر \lr{$T_1$} ثانیه (۲ پالس ساعت) نیز ده ریال دیگر کسر می‌گردد. با رسیدن موجودی به صفر، چراغ \en{Alarm} روشن می‌شود؛ اگر ظرف \lr{$T_2$} ثانیه (۳ پالس ساعت) پس از روشن‌شدن هشدار سکه‌ای اضافه نشود، تماس قطع می‌شود (نشانگر \en{Call} خاموش و \en{Alarm} روشن باقی می‌ماند)؛ در غیر این صورت، با افزایش سکه‌ها تماس ادامه یافته و \en{Alarm} خاموش می‌شود.

\subsection{بررسی و انتخاب روش طراحی}

برای شمارنده‌ی سکه‌ها از دو تراشه‌ی \lr{74192} (شمارنده‌ی بی‌سی‌دی بالارونده/پایین‌رونده) به همراه دو نمایشگر هفت‌قطعه استفاده شد. برای پیاده‌سازی زمان‌سنج \lr{$T_1$} (شمارش ۲ پالس ساعت) از یک فلیپ‌فلاپ استفاده شد، و برای پیاده‌سازی زمان‌سنج \lr{$T_2$} (شمارش ۳ پالس ساعت) از یک تراشه‌ی \lr{74193} استفاده شد.

برای بخش اصلی کنترل حالت‌ها، به دلیل پیچیده‌شدن معادلات و تعداد گیت‌های موردنیاز در روش معمولِ کدگذاری باینری، از \textbf{روش \en{one-hot}} استفاده شد؛ در این روش هر حالت با یک فلیپ‌فلاپ مجزا (در مجموع شش فلیپ‌فلاپ نوع \en{D} از تراشه‌ی \lr{7474}) نمایش داده می‌شود که معادلات حالت بعد و خروجی را به‌شدت ساده می‌کند، هرچند به تعداد فلیپ‌فلاپ بیشتری نسبت به کدگذاری باینری نیاز دارد. با توجه به تعداد کم حالت‌ها (شش حالت)، این افزایش تعداد فلیپ‌فلاپ قابل‌قبول و مقرون‌به‌صرفه‌تر از پیچیدگی معادلات بود.

\subsection{بلوک‌دیاگرام طرح پیشنهادی اولیه}

مدار به‌طور کلی از دو بخش مجزا تشکیل شده است: \textbf{بخش شمارنده‌ی سکه‌ها} (دو تراشه‌ی \lr{74192} به همراه نمایشگرهای هفت‌قطعه) و \textbf{بخش اصلی کنترل تماس} (رجیستر حالت \en{one-hot} به همراه زمان‌سنج‌های \lr{$T_1$} و \lr{$T_2$} و منطق تولید \en{Call} و \en{Alarm}). ارتباط این دو بخش تنها از طریق دو سیگنال برقرار است: سیگنال \en{Deduction} (تولیدشده در بخش اصلی و ارسالی به بخش شمارنده برای کسر یک واحد از موجودی) و سیگنال \en{Z} (تولیدشده در بخش شمارنده و نشان‌دهنده‌ی صفر بودن موجودی سکه‌ها، که به‌عنوان یکی از مهم‌ترین ورودی‌های بخش اصلی به کار می‌رود).

\begin{figure}[H]
\centering
\includegraphics[width=0.85\linewidth]{Screenshot 2026-08-22 082703.png}
\caption{بلوک‌دیاگرام اولیه‌ی مدار تلفن راه‌دور}
\end{figure}

\begin{figure}[H]
\centering
\includegraphics[width=0.85\linewidth]{Screenshot 2026-08-22 082803.png}
\caption{جدول حالات و خروجی هر حالت (این تصویر برای شفافیت و دقت بیشتر بر اساس حالت‌ها و توضیحات ارائه‌شده با کمک هوش مصنوعی بازسازی شده؛ طراحی اصلی مدار توسط خود ما انجام شده است)}
\end{figure}

\subsection{روند پیاده‌سازی}

\paragraph{بخش اول: سیستم سکه‌ها.}
در این بخش از دو تراشه‌ی \lr{74192} (شمارنده‌ی بی‌سی‌دی) استفاده شد که خروجی آن‌ها به نمایشگرهای هفت‌قطعه متصل است؛ ارتباط منطقی این بخش با بخش اصلی مدار تنها از طریق یک ورودی (\en{Deduction}) و یک خروجی (\en{Z}) برقرار می‌شود. سیگنال \en{Deduction} به ورودی \en{Clock Down} این تراشه‌ها متصل است و با فعال‌شدن آن، یک واحد از تعداد سکه‌ها کاسته می‌شود؛ خروجی \en{Z} نیز نشان‌دهنده‌ی صفر شدن موجودی سکه‌هاست و به‌عنوان یکی از مهم‌ترین ورودی‌های بخش دیگر مدار عمل می‌کند.

\begin{figure}[H]
\centering
\includegraphics[width=0.85\linewidth]{Screenshot 2026-08-22 083400.png}
\caption{پیاده‌سازی سیستم شمارنده‌ی سکه‌ها}
\end{figure}

\paragraph{بخش دوم: سیستم اصلی تماس.}
همان‌طور که گفته شد، برای این بخش از روش \en{one-hot} استفاده شده و شش فلیپ‌فلاپ \lr{7474} به کار رفته که ورودی هر یک بر اساس معادلات به‌دست‌آمده از نمودار \en{ASM} زیر محاسبه و سیم‌کشی شده است. حالت‌های در نظر گرفته‌شده عبارت‌اند از: \lr{$S_0$} (آماده به کار، پیش از شروع تماس)، \lr{$S_1$} (لحظه‌ی شروع تماس و کسر اولین ده‌ریال)، \lr{$S_2$} (تماس برقرار، در انتظار سپری‌شدن \lr{$T_1$})، \lr{$S_3$} (پس از \lr{$T_1$}، کسر ده‌ریال بعدی و بازگشت به \lr{$S_2$})، \lr{$S_4$} (موجودی صفر شده، \en{Alarm} روشن، تماس هنوز برقرار، در انتظار سپری‌شدن \lr{$T_2$}) و \lr{$S_5$} (پایان \lr{$T_2$} بدون افزودن سکه، قطع تماس با باقی‌ماندن \en{Alarm}).

\begin{figure}[H]
\centering
\includegraphics[width=0.85\linewidth]{Screenshot 2026-08-22 083616.png}
\caption{نمودار \lr{ASM} بخش اصلی کنترل تماس (این تصویر نیز برای شفافیت و دقت بیشتر بر اساس حالت‌ها و توضیحات ارائه‌شده با کمک هوش مصنوعی بازسازی شده؛ طراحی اصلی مدار توسط خود ما انجام شده است)}
\end{figure}

\paragraph{زمان‌سنج‌های \lr{$T_1$} و \lr{$T_2$}.}
برای شمارش دو پالس ساعت لازم برای \lr{$T_1$} از یک فلیپ‌فلاپ استفاده شد و برای شمارش سه پالس ساعت لازم برای \lr{$T_2$} از یک تراشه‌ی \lr{74193} استفاده شد. ورودی \en{Master Reset} تراشه‌ی مربوط به \lr{$T_2$} به‌گونه‌ای به \lr{$\overline{S_4}$} (متمم سیگنال حالت \lr{$S_4$}) متصل شده که این شمارنده تنها در حالت \lr{$S_4$} فعال باشد و با خروج از این حالت، مقدار آن بازنشانی شود.

\begin{figure}[H]
\centering
\includegraphics[width=0.85\linewidth]{Screenshot 2026-08-22 083707.png}
\caption{پیاده‌سازی زمان‌سنج \lr{$T_1$} با فلیپ‌فلاپ و زمان‌سنج \lr{$T_2$} با تراشه‌ی \lr{74193}}
\end{figure}

\paragraph{مقداردهی اولیه‌ی خودکار.}
برای آنکه در همان لحظه‌ی شروع شبیه‌سازی --- بدون نیاز به فشردن دستی کلید \en{Reset} --- مدار به‌طور خودکار در حالت \lr{$S_0$} قرار گیرد، از یک تراشه‌ی نظارت بر تغذیه/بازنشانی (\lr{MCP809-300}) استفاده شد که با روشن‌شدن تغذیه، یک پالس بازنشانی کوتاه تولید کرده و رجیستر \en{one-hot} را در حالت \lr{$S_0$} مقداردهی می‌کند.

\begin{figure}[H]
\centering
\includegraphics[width=0.85\linewidth]{Screenshot 2026-08-22 083826.png}
\caption{مدار مقداردهی اولیه‌ی خودکار با تراشه‌ی \lr{MCP809-300}}
\end{figure}

\paragraph{سیستم چراغ‌ها و سیگنال \lr{Deduction}.}
سیگنال \en{Deduction} پیش از اتصال به ورودی \en{Clock Down} شمارنده‌ی سکه‌ها، با متمم سیگنال \en{Z} در یک گیت \en{NAND} ترکیب می‌شود؛ دلیل این کار آن است که ورودی \en{Clock Down} باید همواره در سطح منطقی ۱ بماند مگر در لحظه‌ای که کسر واقعی باید رخ دهد، تا از یک‌سو عملکرد شمارش در هر مرحله به‌درستی انجام شود و از سوی دیگر، پس از رسیدن موجودی به صفر، شمارنده به مقادیر ۹۹ یا کمتر (سرریز به عقب) بازنگردد.

\begin{figure}[H]
\centering
\includegraphics[width=0.85\linewidth]{Screenshot 2026-08-22 084012.png}
\caption{بخشی از منطق تولید سیگنال‌های \lr{Call} و \lr{Alarm}}
\end{figure}

\begin{figure}[H]
\centering
\includegraphics[width=0.85\linewidth]{Screenshot 2026-08-22 084427.png}
\caption{سیستم نهایی چراغ‌های \lr{LED} و منطق سیگنال \lr{Deduction}}
\end{figure}

\subsection{دیاگرام/شماتیک طرح اولیه}

طرح اولیه‌ی این سیستم تفاوت محسوسی با طرح نهایی نداشت؛ تنها تفاوت آن بود که در طرح اولیه، مقداردهی اولیه‌ی رجیستر حالت به حالت \lr{$S_0$} صرفاً از طریق فشردن کلید \en{Reset} توسط کاربر انجام می‌شد. در طرح نهایی، مدار \lr{MCP809-300} برای انجام خودکار همین کار در لحظه‌ی روشن‌شدن تغذیه به مدار اضافه شد تا نیاز به دخالت دستی در ابتدای هر اجرای شبیه‌سازی از بین برود.

\subsection{دیاگرام طرح نهایی}

شماتیک کامل و نهایی مدار تلفن راه‌دور (خروجی مستقیم از \en{Proteus}) در ادامه آمده است. در سمت چپ این شماتیک، بخش شمارنده‌ی سکه‌ها (تراشه‌های \lr{74192} و نمایشگرهای هفت‌قطعه) و در سمت راست، رجیستر حالت \en{one-hot} (شش فلیپ‌فلاپ \lr{7474})، زمان‌سنج‌های \lr{$T_1$} و \lr{$T_2$}، مدار مقداردهی اولیه‌ی \lr{MCP809-300} و منطق نهایی چراغ‌های \en{Call}/\en{Alarm} قابل مشاهده است.



\begin{figure}[H]
\centering
\includegraphics[width=\linewidth]{4-2-main.png}
\caption{شماتیک نهایی مدار تلفن راه‌دور در \lr{Proteus}}
\end{figure}

\subsection{بررسی \lr{Test Case} ها}

\paragraph{الف) برقراری یک تماس عادی.} ابتدا با فشردن کلید \en{Coin} چند سکه به موجودی اضافه و سپس \en{Start} فشرده شد. بلافاصله ده ریال کسر و چراغ سبز \en{Call} روشن شد؛ سپس در هر \lr{$T_1$} (۲ پالس ساعت)، مدار یک بار دیگر ده ریال کسر کرده و به حالت انتظار برای دوره‌ی بعد بازگشت.

\begin{figure}[H]
\centering
\includegraphics[width=0.75\linewidth]{4-2-1 idle state.png}
\caption{حالت اولیه (\lr{$S_0$}) پیش از انداختن سکه}
\end{figure}

\begin{figure}[H]
\centering
\includegraphics[width=0.75\linewidth]{4-2-2 adding coin.png}
\caption{افزایش موجودی با فشردن کلید \lr{Coin}}
\end{figure}

\begin{figure}[H]
\centering
\includegraphics[width=0.75\linewidth]{4-2-3 state s1 and coin initial deduction after the clock.png}
\caption{ورود به \lr{$S_1$} و کسر اولین ده‌ریال پس از پالس ساعت}
\end{figure}

\begin{figure}[H]
\centering
\includegraphics[width=0.75\linewidth]{4-2-4 entering state s2 after initial deduction in s1 .png}
\caption{ورود به \lr{$S_2$} پس از کسر اولیه در \lr{$S_1$}}
\end{figure}

\begin{figure}[H]
\centering
\includegraphics[width=0.75\linewidth]{4-2-5 entering state s3 after 2 clock pulse.png}
\caption{ورود به \lr{$S_3$} پس از سپری‌شدن \lr{$T_1$} (۲ پالس ساعت)}
\end{figure}

\begin{figure}[H]
\centering
\includegraphics[width=0.75\linewidth]{4-2-6 deduction and going back to s2.png}
\caption{کسر ده‌ریال بعدی در \lr{$S_3$} و بازگشت به \lr{$S_2$}}
\end{figure}

\begin{figure}[H]
\centering
\includegraphics[width=0.75\linewidth]{4-2-7 green led on when call is going on.png}
\caption{روشن ماندن چراغ سبز \lr{Call} در طول تماس}
\end{figure}

\paragraph{ب) اتمام موجودی و قطع خودکار تماس.} با ادامه‌ی تماس تا صفر شدن موجودی سکه‌ها، مدار وارد \lr{$S_4$} شده و چراغ قرمز \en{Alarm} روشن می‌شود درحالی‌که تماس هنوز برقرار است (هر دو چراغ روشن). چون در این آزمون ظرف \lr{$T_2$} (۳ پالس ساعت) سکه‌ای اضافه نشد، طبق انتظار تماس قطع (چراغ سبز خاموش) و چراغ قرمز روشن باقی ماند (\lr{$S_5$}).

\begin{figure}[H]
\centering
\includegraphics[width=0.75\linewidth]{4-2-8 entering s4 after getting zero on coins.png}
\caption{ورود به \lr{$S_4$} پس از صفر شدن موجودی سکه‌ها}
\end{figure}

\begin{figure}[H]
\centering
\includegraphics[width=0.75\linewidth]{4-2-9 both leds on when s4.png}
\caption{روشن بودن هم‌زمان هر دو چراغ \lr{Call} و \lr{Alarm} در \lr{$S_4$}}
\end{figure}

\begin{figure}[H]
\centering
\includegraphics[width=0.75\linewidth]{4-2-10 call ended and led red still glowing and state s5.png}
\caption{قطع تماس پس از پایان \lr{$T_2$} --- چراغ قرمز روشن باقی می‌ماند (\lr{$S_5$})}
\end{figure}

\paragraph{ج) افزودن سکه در حین هشدار و ادامه‌ی تماس.} در آزمون دیگری، هنگامی‌که مدار در حالت هشدار (\lr{$S_4$}) قرار داشت، پیش از پایان \lr{$T_2$} یک سکه اضافه شد. طبق انتظار، چراغ \en{Alarm} خاموش شده و تماس بدون قطعی ادامه یافت.

\begin{figure}[H]
\centering
\includegraphics[width=0.75\linewidth]{4-2-11 enterning s4 in another call.png}
\caption{ورود مجدد به حالت هشدار در یک تماس دیگر}
\end{figure}

\begin{figure}[H]
\centering
\includegraphics[width=0.75\linewidth]{4-2-12 adding coins and continue the call.png}
\caption{افزودن سکه در حین هشدار و ادامه‌ی تماس بدون قطعی}
\end{figure}

\subsection{نتیجه‌گیری بخش دوم}

مدار کنترل‌کننده‌ی تلفن راه‌دور با ترکیب یک شمارنده‌ی بی‌سی‌دی برای موجودی سکه‌ها و یک رجیستر حالت \en{one-hot} شش‌حالته برای منطق تماس، به‌درستی طراحی و پیاده‌سازی شد. هر سه سناریوی آزمایشی --- برقراری تماس عادی، اتمام موجودی و قطع خودکار تماس، و افزودن سکه در حین هشدار --- رفتار مطابق با صورت مسئله را نشان دادند. انتخاب روش \en{one-hot} برای این بخش، با وجود افزایش تعداد فلیپ‌فلاپ‌ها، معادلات منطقی را به‌طور محسوسی ساده کرد و اشکال‌زدایی و اطمینان از صحت هر انتقال حالت را آسان‌تر ساخت.


\section{نتیجه‌گیری کلی}


در این آزمایش، روند کامل طراحی یک مدار کنترل‌کننده --- از تعریف مسئله و ترسیم نمودار \en{ASM} تا پیاده‌سازی گیتی و صحت‌سنجی عملی --- برای دو مسئله‌ی متفاوت تمرین شد. مقایسه‌ی این دو پیاده‌سازی نشان داد که انتخاب روش کدگذاری حالت (باینری همراه با دیکودر، یا \en{one-hot}) باید بر اساس تعداد حالت‌ها و پیچیدگی معادلات حالت بعد/خروجی انجام شود: برای تعداد حالت کم با گذارهای نسبتاً پیچیده (تلفن راه‌دور) روش \en{one-hot} ساده‌تر و قابل‌اطمینان‌تر بود، درحالی‌که برای توالی‌های خطی‌تر (تایمر ماشین لباس‌شویی) استفاده از رجیستر باینری به‌همراه دیکودر مصرف قطعات را بدون افزایش پیچیدگی معادلات کاهش داد. در هر دو مدار، تمام سناریوهای آزمایشی تعریف‌شده با موفقیت اجرا شدند و رفتار مدار در تمام حالت‌های مرزی (اتمام زمان هر مرحله، فشردن \en{Reset}، صفر شدن موجودی سکه‌ها) مطابق با صورت مسئله بود.

\end{document}